% !TeX root = ../../main.tex
% ------------------------------------------------------------------
%  4.2.4 Conexiones, puntos de falla y contingencia
%  Fragmento del Subdocumento 4.2, traido desde contenido.tex con
%  \parte{partes/02_c_conexiones}. Las rutas son relativas a la carpeta del
%  subdocumento, nunca a la de main.tex.
%
%  Coherente con la arquitectura de despliegue (puntos únicos de falla,
%  RT-02.11) y con la autonomía declarada en 4.2.1 (Tabla t29).
% ------------------------------------------------------------------

\section{Conexiones, puntos de falla y contingencia}
\label{sec:conexiones}

Esta sección describe cómo se conectan los sitios, el terreno y la nube, identifica cada punto donde esa conexión o su infraestructura puede fallar, y declara cómo se resuelve cada falla o, cuando no se resuelve de forma automática, cuál es la contingencia. Su punto de partida es el caso: la conectividad actual de Puelche se corta con frecuencia y en algunos lugares no existe, por lo que la solución no puede suponer un enlace disponible y debe decir, para cada conexión, qué pasa cuando no lo está.

\subsection{Conexiones entre sitios, terreno y nube}

Los dos dominios fijos, la nube y el on-premise, se conectan por túneles VPN IPsec que terminan en el firewall de cada sitio (D-01) y en el Transit Gateway de AWS en sa-east-1, con enrutamiento dinámico BGP; el mismo Transit Gateway enruta entre los sitios, de modo que Concepción y los cross-docking alcanzan Talca sin exponerla a Internet. Cada sitio tiene dos caminos físicamente independientes: fibra y LTE en los centros de distribución, y Starlink y LTE de dos proveedores en los cross-docking, que no tienen fibra. El terreno no pasa por los sitios: los terminales de preventa y de reparto hablan directamente con API Gateway por la red celular del operador. La Figura~\ref{fig:conexiones} muestra esas conexiones.

\begin{figure}[htbp]
  \centering
  \LFXcuerpofigura
  \begin{tikzpicture}[
      nodo/.style={draw=lafroxGris, line width=0.5pt, rounded corners=1.5pt,
                   inner sep=4pt, align=center},
      sitio/.style={nodo, text width=2.55cm, minimum height=1.25cm},
      flecha/.style={-{latex}, line width=0.6pt, color=lafroxTinta},
      interna/.style={-{latex}, line width=0.6pt, color=lafroxGris, dashed},
      rotulo/.style={fill=white, inner sep=1.5pt, align=center}]
    % Usuarios externos
    \node[nodo, text width=3.9cm] (usr) at (10.7,0.45)
      {Usuarios externos: clientes, transportistas y proveedores};
    % Región primaria
    \draw[lafroxTinta, line width=0.6pt, rounded corners=2pt] (0,-1.3) rectangle (12.2,-3.4);
    \node[anchor=north west, font=\intersemibold, inner sep=0pt] at (0.15,-1.45) {AWS sa-east-1};
    \node[nodo, text width=2.2cm] (vgw)  at (1.75,-2.6) {Transit Gateway};
    \node[nodo, text width=2.6cm] (vpce) at (4.95,-2.6) {VPC Endpoints de SQS, IoT Core, SSM y ECR};
    \node[nodo, text width=2.0cm] (apig) at (8.0,-2.6) {API Gateway};
    \node[nodo, text width=2.0cm] (cf)   at (10.7,-2.6) {CloudFront y WAF};
    % Región secundaria
    \draw[lafroxTinta, line width=0.6pt, rounded corners=2pt] (12.75,-1.3) rectangle (15.7,-3.4);
    \node[align=center] (us) at (14.22,-2.35) {\intersemibold AWS us-east-1\\recuperación\\y réplica};
    \draw[flecha] (12.2,-2.35) -- (12.75,-2.35);
    % Sitios y terreno
    \node[sitio] (tal) at (1.45,-6.4)  {\intersemibold CD Talca\\\normalfont D-01 en par};
    \node[sitio] (con) at (4.55,-6.4)  {\intersemibold CD Concepción\\\normalfont D-01 de borde};
    \node[sitio] (cdk) at (7.65,-6.4)  {\intersemibold Cross-docking\\\normalfont E-01, tres sitios};
    \node[sitio] (ter) at (10.75,-6.4) {\intersemibold Terreno\\\normalfont C-01 y C-02};
    \node[sitio] (cam) at (13.85,-6.4) {\intersemibold Camiones\\\normalfont B-03};
    % Enlaces hacia la nube
    \draw[flecha] (usr) -- node[rotulo, left=3pt, pos=0.3] {HTTPS} (cf);
    \draw[flecha] (tal.north) -- node[rotulo, pos=0.5] {fibra D-03\\LTE D-04\\IPsec} (vgw.south);
    \draw[flecha] (con.north) -- node[rotulo, pos=0.5] {fibra D-03\\LTE D-04\\IPsec} (vgw.south east);
    \draw[flecha] (cdk.north) -- node[rotulo, pos=0.35] {Starlink D-06\\LTE D-04\\IPsec} (vgw.east);
    \draw[flecha] (ter.north) -- node[rotulo, pos=0.5] {red celular\\HTTPS} (apig.south);
    % Sincronizaciones y descargas locales
    \draw[interna] (con.south) -- ++(0,-0.4) -| node[rotulo, below=1pt, pos=0.25] {TCP 8080 y 5432} (tal.south);
    \draw[interna] (cdk.south) -- ++(0,-0.95) -| node[rotulo, below=1pt, pos=0.25] {AMQPS 5671 entre brokers} ([xshift=0.45cm]tal.south west);
    \draw[interna] (cam.south) -- ++(0,-0.4) -| node[rotulo, below=1pt, pos=0.25] {BLE con el terminal\\del conductor} (ter.south);
  \end{tikzpicture}
  \caption{Conexiones entre los sitios, el terreno y la nube}
  \label{fig:conexiones}
\end{figure}

La figura distingue dos tipos de conexión. Las líneas continuas son conexiones hacia la nube, y todas se inician desde el sitio o desde el dispositivo, sin puertos abiertos hacia el on-premise. Las únicas excepciones son la réplica de AWS DMS, que lee el registro de escritura de PostgreSQL en VM-02, y el trabajador Laravel \texttt{erp-sync}, que consulta la capa anticorrupción en VM-04; ambas entran por el túnel IPsec autenticado y quedan auditadas. Las líneas discontinuas son sincronizaciones entre sitios o descargas locales, que no dependen de la nube: Concepción sincroniza con Talca por la VPN, los cross-docking envían a Talca el detalle de su WMS entre brokers por la VPN, y los termógrafos sincronizan por Bluetooth con el terminal del conductor durante la ruta —que alerta la excursión térmica en el momento y la reenvía a la nube cuando recupera cobertura— y, cuando el camión vuelve, le descargan el registro completo, que el terminal envía a la nube. Los eventos críticos de los cross-docking van directo a SQS FIFO, de modo que la operación de esas plataformas no depende de Talca.

La Tabla~\ref{tab:conexiones} resume cada conexión con su medio, su camino alternativo y el tiempo de conmutación entre ambos.

\begin{tablalafrox}{Conexiones de la solución y su camino alternativo}{tab:conexiones}%
  {F{0.24} Y F{0.24} F{0.16}}%
  {\cab{Conexión} & \cab{Medio y protocolo} & \cab{Camino alternativo} & \cab{Conmutación}}
  CD Talca con la nube & Fibra D-03, túneles IPsec con BGP & LTE D-04 & Menos de 30 s \\
  CD Concepción con la nube & Fibra D-03, túneles IPsec con BGP & LTE D-04 & Menos de 30 s \\
  Cross-docking con la nube & Starlink D-06, túneles IPsec con BGP & LTE D-04 de dos proveedores & Menos de 30 s \\
  Concepción con Talca & TCP 8080 y 5432 con TLS sobre la VPN & Operación autónoma & No aplica \\
  Cross-docking con Talca & AMQPS 5671 entre brokers, por la VPN & Envío diferido & No aplica \\
  Terreno con la nube & Red celular del operador, HTTPS & Almacén local del dispositivo & No aplica \\
  Termógrafos con el terminal del conductor & BLE con el terminal del conductor, en ruta y al regresar & Registro interno del termógrafo & No aplica \\
  Usuarios externos con los portales & Internet, HTTPS por CloudFront & Ninguno: el autoservicio requiere conexión & No aplica \\
  Región primaria con la secundaria & Replicación nativa de Aurora, DynamoDB y S3 & Respaldo inmutable & Promoción autorizada \\
\end{tablalafrox}

Solo tres conexiones tienen un segundo camino de red, y son las que unen un sitio con la nube. Las demás no lo necesitan porque su alternativa no es otro enlace sino la operación sin conexión: Concepción opera sin Talca, el terreno opera sin señal y el termógrafo guarda su propio registro. Por eso la redundancia de enlaces reduce la frecuencia de las desconexiones, pero no es condición para operar. El portal de clientes es la única conexión sin alternativa, y eso es deliberado: el autoservicio no se ofrece sin conexión.

\subsection{Puntos de falla en los sitios y en los enlaces}

Cada conexión se apoya en equipos del sitio que también pueden fallar. La Tabla~\ref{tab:fallas-sitios} recorre esos puntos de falla, desde el enlace hasta la energía del recinto, e indica para cada uno el mecanismo que la resuelve y la contingencia si ese mecanismo no basta.

\begin{tablalafrox}{Puntos de falla en los sitios y en los enlaces}{tab:fallas-sitios}%
  {F{0.24} Y Y}%
  {\cab{Punto de falla} & \cab{Resolución} & \cab{Contingencia}}
  Fibra de un centro de distribución & BGP conmuta los túneles al LTE en menos de 30 s. & Si cae también el LTE, operación local 24 h. \\
  Starlink de un cross-docking & El LTE de dos proveedores toma el tráfico en menos de 30 s. & La ventana de 3 h opera 100\,\% local y sincroniza al reconectar. \\
  Enlace de Los Ángeles en la madrugada & Starlink queda como camino único, porque la red móvil falla a esa hora. & La ventana de madrugada opera íntegramente local. \\
  Firewall de Talca & La unidad pasiva asume las direcciones y los túneles en menos de 30 s. & Operación local 24 h sin enlace. \\
  Firewall de Concepción & Equipo único; se repone. & Concepción opera 24 h sin enlace mientras se reemplaza. \\
  Nodo del clúster de Talca & Las VM se reinician en los nodos restantes, con sus datos replicados por Ceph. & Ante un segundo nodo, el DRP local promueve el WMS de Concepción en 1 a 2 h. \\
  Instancia de escritura VM-02 & Se reinicia en otro nodo del clúster. & D-05 restaura la base en hasta 4 h sin enlace. \\
  Disco de un nodo de Talca & El RAID 10 reconstruye sobre el repuesto en caliente. & Ceph conserva una segunda copia en otro nodo. \\
  Miembro del stack de switches de Talca & El otro miembro mantiene el tráfico. & Reposición sin corte. \\
  Servidor de borde de Concepción & El RAID 10 tolera la falla de un disco; el equipo tiene fuente única y, ante su falla, se repone. & La pila se restaura desde la réplica en Aurora. \\
  Mini-PC de un cross-docking & Equipo único; se repone. & El nuevo equipo se resincroniza desde Talca y SQS FIFO. \\
  Firewall de un cross-docking & Equipo único; se repone. & La ventana de 3 h opera 100\,\% local y sincroniza al reconectar. \\
  Periféricos de andén & Sin redundancia por equipo. & La etiqueta se emite en otra de las impresoras del centro de distribución. \\
  Energía de la sala de Talca & UPS N+1 de 30 min y generador que toma carga en 8 a 15 s. & Estanque de 24 h y contrato de reabastecimiento. \\
  Climatización de la sala de Talca & Segunda unidad de precisión en N+1. & Alerta del monitoreo de los sensores ambientales. \\
\end{tablalafrox}

La tabla muestra que los puntos únicos de falla que quedan están donde la redundancia no compensa su costo: el firewall y el servidor de Concepción, el firewall y el mini-PC de cada cross-docking y los periféricos de andén. En los cuatro primeros casos la contingencia es la misma autonomía que ya existe para los cortes de enlace. El firewall y el servidor de Concepción y el firewall, el mini-PC y el switch de cada cross-docking operan, además, con fuente única; esa excepción a la exigencia de fuente redundante se declara con el mismo fundamento. En la ruta de la ventana de despacho de Talca quedan dos elementos sin duplicar: la instancia de escritura VM-02, que se reinicia en otro nodo del clúster, y cada impresora de andén, cuya etiqueta puede emitir otra de las impresoras del centro de distribución; el resto de esa ruta tiene redundancia en cada capa. Los mecanismos de alta disponibilidad que sostienen estas resoluciones se desarrollan en la sección de arquitectura de despliegue.

\subsection{Puntos de falla en la nube, el terreno y las integraciones}

La Tabla~\ref{tab:fallas-nube} completa el recorrido con los puntos de falla que están fuera de los sitios: la nube, el terreno y los sistemas de terceros.

\begin{tablalafrox}{Puntos de falla en la nube, el terreno y las integraciones}{tab:fallas-nube}%
  {F{0.24} Y Y}%
  {\cab{Punto de falla} & \cab{Resolución} & \cab{Contingencia}}
  Zona de disponibilidad de sa-east-1 & Aurora conmuta en menos de 30 s, ElastiCache en menos de 60 s y Fargate reprograma las tareas. & La operación continúa en las zonas restantes. \\
  Región sa-east-1 completa & Route 53 redirige el tráfico y, con autorización del CLIENTE, se promueve la réplica en us-east-1. & RTO de 4 h y RPO de 15 min; la bodega y el terreno operan en local, y los terminales de bodega trabajan contra la puerta de API local del WMS. \\
  Túnel VPN durante la replicación & DMS encola los cambios en el origen. & Reanuda sin pérdida al reconectar; la réplica queda desactualizada durante el corte. \\
  IdP maestro o enlace de identidad & La caché de solo lectura, de 24 h, conserva los datos de identidad recibidos, y el verificador local habilita los relevos de turno con el manifiesto firmado y el PIN personal. & La credencial de turno dura hasta 8 h en bodega y 14 h en terreno; no hay revocación remota mientras dure el corte; cuenta de emergencia sellada. \\
  Señal móvil en ruta & La aplicación opera contra el almacén local del dispositivo. & Sincroniza al reconectar dentro de los 10 min de la Tabla~\ref{tab:t29}. \\
  ERP legado & El trabajador \texttt{erp-sync} retiene las operaciones en su cola SQS hasta 24 h y reintenta con la misma clave idempotente, de modo que un reintento no emite otro documento. & El camión no se libera sin guía confirmada o contingencia tributaria aprobada por el CLIENTE. \\
  Perfil de API saturado & El balanceador reparte entre tareas y el escalado agrega tareas cuando los procesos PHP-FPM ocupados superan el 70 \%. & Limitación de tasa en API Gateway con mensaje explícito al usuario. \\
  Mensaje de reconciliación inválido o que falla & El consumidor rechaza el sobre de versión desconocida sin aplicarlo; tras cinco intentos pasa a la cola de mensajes fallidos. & Alarma inmediata y reproceso manual trazable; el grupo afectado no bloquea a los demás. \\
  Excursión térmica sin enlace & El gateway bloquea el despacho en menos de 5 s en el borde. & Buffer de 24 h y alerta al reconectar. \\
  SII, cadenas o notificaciones & Reintento con espera creciente y cola de mensajes fallidos. & Reproceso desde la cola al recuperarse el tercero. \\
  Carga del peak de septiembre & Escalado automático e independiente de los perfiles de API y de trabajadores en Fargate, y de Aurora. & Limitación de tasa en API Gateway. \\
\end{tablalafrox}

Ninguna de estas fallas detiene la bodega ni el terreno. La peor de ellas, la pérdida de toda la región primaria, deja sin servicio a los portales y a la analítica hasta la promoción de la réplica, pero la preparación, el despacho y la entrega siguen operando contra sus bases locales y sus dispositivos, porque ninguna de esas transacciones depende de la nube para confirmarse.
